\medskip

\begin{enumerate}
\item $x$ représente les membres présents le lundi et absents le jeudi ; 

Avec la ligne 1 : comme $45 + x = 75$, on a donc $x = 75 - 45 = 30$.

Avec la colonne 2 : comme $x + 5 = 35$, on a donc $x = 35 - 5 = 30$.`
\item
	\begin{enumerate}
		\item 5 adhérents ne sont venus ni le lundi ni le jeudi ; la probabilité est égale à $\dfrac{5}{100} = 0,05$.
		\item Il y a 30 personnes présentes le lundi et absentes le jeudi et 20 personnes présentes le jeudi (donc différentes ds 30 personnes précédentes) et absente le lundi, soit en tout 50 personnes venues un seul jour ; la probabilité est égale à $\dfrac{50}{100} = 0,5$.
		\item Sur les 75 adhérents venus le lundi, 45 sont venus le jeudi ; la probabilité est égale à $\dfrac{45}{75} = \dfrac{3 \times 15}{3 \times 25} = \dfrac{15}{25} = \dfrac{15 \times 4}{25 \times 4} = \dfrac{60}{100} = 0,6$.
	\end{enumerate}
\item 
	\begin{enumerate}
		\item 100 adhérents versent 100~\euro{} soit en tout $100 \times 100 = \np{10000}$~(\euro).
		\item De 2026 à 2021 il y aura 15 augmentations. Chaque année il y aura 5 adhérents de plus donc $5 \times 100 = 500$~\euro versés en plus au club.
		
La somme totale versée au club sera :

$\np{10000} + \np{10000} + 500 + \np{10000} + 500 + \ldots + \np{10000} + 15 \times 500 = \dfrac{2 \times \np{10000} + 15 \times 500}{2} \times (15 + 1) = 8 \times (\np{20000} + \np{7500}) = 8 \times \np{27500} = \np{220000}$~(\euro).
	\end{enumerate}
\end{enumerate}
